\documentclass[11pt,a4paper]{article}
\usepackage[utf8]{inputenc}
\usepackage[margin=1in]{geometry}
\usepackage{xcolor}
\usepackage{tcolorbox}
\usepackage{hyperref}
\usepackage{listings}
\usepackage{enumitem}
\usepackage{booktabs}

% Define elegant, modern color scheme
\definecolor{primary}{HTML}{1E293B}   % Deep Slate (Navy/Charcoal)
\definecolor{secondary}{HTML}{0F766E} % Rich Teal
\definecolor{accent}{HTML}{B45309}    % Amber / Orange
\definecolor{success}{HTML}{15803D}   % Forest Green
\definecolor{background}{HTML}{F8FAFC} % Slate 50 (Very light gray-blue)
\definecolor{codebg}{HTML}{F1F5F9}     % Slate 100

% Configure Hyperref
\hypersetup{
    colorlinks=true,
    linkcolor=secondary,
    filecolor=secondary,      
    urlcolor=secondary,
    pdftitle={AI Animation Platform - Documentation},
    pdfauthor={Antigravity AI Pair Programming},
}

% Configure custom tcolorboxes
\newtcolorbox{setupbox}[2][]{
    colback=background,
    colframe=secondary,
    fonttitle=\bfseries,
    title={#2},
    coltitle=white,
    arc=2mm,
    #1
}

\newtcolorbox{guardrailbox}[2][]{
    colback=background,
    colframe=accent,
    fonttitle=\bfseries,
    title={#2},
    coltitle=white,
    arc=2mm,
    #1
}

\newtcolorbox{bugbox}[2][]{
    colback=background,
    colframe=primary,
    fonttitle=\bfseries,
    title={#2},
    coltitle=white,
    arc=2mm,
    #1
}

% Configure Listings for source code display
\lstset{
    backgroundcolor=\color{codebg},
    commentstyle=\color{gray},
    keywordstyle=\color{secondary}\bfseries,
    numberstyle=\tiny\color{gray},
    stringstyle=\color{accent},
    basicstyle=\ttfamily\small,
    breakatwhitespace=false,         
    breaklines=true,                 
    captionpos=b,                    
    keepspaces=true,                 
    numbers=left,                    
    numbersep=5pt,                  
    showspaces=false,                
    showstringspaces=false,
    showtabs=false,                  
    tabsize=4,
    language=Python
}

\title{\Huge \textbf{AI Educational Animation Platform} \\ \Large Dynamic Manim Scene Generation \& Automation Pipeline}
\author{\textbf{System Documentation \& Setup Guide}}
\date{May 30, 2026}

\begin{document}

\maketitle
\tableofcontents
\newpage

\section{Introduction \& Project Goal}
The \textbf{AI Educational Animation Platform} is a Python-based pipeline that automatically compiles high-quality mathematical and educational animations from human-readable text prompts. It connects the power of large language models (specifically \textbf{Gemini 2.5 Flash} using the modern \texttt{google-genai} SDK) to \textbf{Manim (Community Edition)}, a math animation engine.

The core objective is to allow educators or developers to type a concept (such as \textit{"explain binary search visually"}) and obtain a compiled, ready-to-preview MP4 video file in seconds.

---

\section{System Environment \& Setup}
To run this pipeline on Windows, we initialize a dedicated virtual environment to handle Python dependencies and the underlying LaTeX/FFmpeg tools required by Manim.

\begin{setupbox}{Step 1: Create and Activate the Virtual Environment}
First, open your PowerShell or Command Prompt in the project workspace folder and initialize the environment:
\begin{lstlisting}[language=bash]
# 1. Create a Python Virtual Environment
python -m venv venv

# 2. Activate the Virtual Environment (PowerShell)
.\venv\Scripts\Activate.ps1

# 3. Activate the Virtual Environment (Command Prompt)
.\venv\Scripts\activate.bat
\end{lstlisting}
\end{setupbox}

\begin{setupbox}{Step 2: Install Python Dependencies}
Next, install the required packages. Ensure your virtual environment is active (indicated by a \texttt{(venv)} prefix in the terminal prompt):
\begin{lstlisting}[language=bash]
pip install -r requirements.txt
\end{lstlisting}
\end{setupbox}

The file \texttt{requirements.txt} contains:
\begin{lstlisting}[language=Python]
google-genai
python-dotenv
manim
\end{lstlisting}

\begin{setupbox}{Step 3: Setup the Environment Variables}
Create a \texttt{.env} file in the root workspace folder containing your API key:
\begin{lstlisting}[language=bash]
GEMINI_API_KEY=your_gemini_api_key_here
\end{lstlisting}
\end{setupbox}

\begin{setupbox}{Step 4: Verify System Dependencies (Manim Core)}
Manim relies on underlying system applications. If you do not have them installed, run:
\begin{itemize}
    \item \textbf{FFmpeg}: Required for video encoding. Install via \texttt{choco install ffmpeg} or scoop.
    \item \textbf{LaTeX (Optional, for math formulas)}: MikTeX or TeX Live is recommended.
\end{itemize}
\end{setupbox}

---

\section{Project File Structure}
The workspace is organized as follows:
\begin{itemize}[label=\(\bullet\)]
    \item \textbf{src/llm\_client.py}: Interface for the Gemini client, including retry logic and system instructions (guardrails).
    \item \textbf{src/pipeline.py}: Coordinates prompt input, triggers LLM generation, extracts Python code, detects the Scene class name, and runs Manim as a subprocess.
    \item \textbf{generated\_scene.py}: The script dynamically written by the pipeline during each run.
    \item \textbf{media/}: Directory where compiled MP4 animations are saved by Manim.
\end{itemize}

---

\section{The Core Automation Pipeline}

\subsection{1. Code Generation (Gemini Bridge)}
The client connects to the Gemini API using the \texttt{gemini-2.5-flash} model. It uses strict system instructions to enforce clean, executable output:

\begin{lstlisting}[language=Python]
# src/llm_client.py excerpt
response = client.models.generate_content(
    model='gemini-2.5-flash',
    contents=user_instruction,
    config=types.GenerateContentConfig(
        system_instruction=system_instruction,
        temperature=0.1,  # Kept deterministic
    )
)
\end{lstlisting}

\subsection{2. Dynamic Scene Parsing}
Once code is returned, \texttt{src/pipeline.py} uses regular expressions to find which class name was generated (e.g. \texttt{BinarySearch}) so Manim knows what class to target for rendering:
\begin{lstlisting}[language=Python]
# src/pipeline.py excerpt
pattern = r"class\s+(\w+)\s*\(\s*\w*Scene\w*\s*\)\s*:"
match = re.search(pattern, code)
\end{lstlisting}

\subsection{3. Subprocess Compilation}
The pipeline invokes the Manim CLI tool to render the animation. It runs in low-quality mode (\texttt{-ql}) for speed and opens the preview file (\texttt{-p}) immediately upon completion:
\begin{lstlisting}[language=Python]
command = ["manim", "-pql", filepath, class_name]
subprocess.run(command, check=True)
\end{lstlisting}

---

\section{Debugging \& Optimization History}

During testing of the \textbf{Binary Search} animation concept, three major visual and architectural bugs were found and resolved:

\begin{bugbox}{Bug 1: Numpy Array ValueError in get\_pointer()}
\textbf{Symptom:} The script crashed during compiling with: \\
\texttt{ValueError: The truth value of an array with more than one element is ambiguous.} \\
\textbf{Cause:} In Manim, direction indicators like \texttt{DOWN} are numpy arrays. The code generated had: \\
\texttt{rotate(PI if direction == DOWN else 0)} \\
Evaluating an equality condition directly on numpy arrays evaluates element-wise, triggering the ambiguity exception. \\
\textbf{Resolution:} Replaced with \texttt{all(direction == DOWN)} or \texttt{np.array\_equal()}.
\end{bugbox}

\begin{bugbox}{Bug 2: Bottom Status and Result Text Collisions}
\textbf{Symptom:} The final message \texttt{"Binary Search Complete!"} overlapped directly on top of the search comparison status texts (e.g. \texttt{"70 == 70! Target Found!"}). \\
\textbf{Cause:} The generated code had both texts positioned at \texttt{next\_to(..., DOWN * 3)}, and the previous status text was never cleared/faded out. \\
\textbf{Resolution:} Moved status texts to the screen edge using \texttt{.to\_edge(DOWN, buff=0.5)} and added an explicit \texttt{FadeOut(comparison\_text)} animation before rendering the final result.
\end{bugbox}

\begin{bugbox}{Bug 3: Horizontal Video Edge Overflow (Cutoff)}
\textbf{Symptom:} The left and rightmost boxes of the sorted array (representing elements \texttt{10} and \texttt{90}) were cut off on the screen margins. \\
\textbf{Cause:} The generated layout used boxes of width \texttt{1.0} with a huge buffer spacing of \texttt{0.8}. This resulted in a layout width of \texttt{15.4} units, exceeding Manim's screen width of \texttt{14.22} units. \\
\textbf{Resolution:} Reduced cell width to \texttt{0.9} and the arrangement buffer to \texttt{0.15}, keeping the total width safely under \texttt{9.3} units.
\end{bugbox}

---

\section{System Instructions \& Guardrails}
To prevent the LLM from making similar layout mistakes in future runs, the system instructions inside \texttt{src/llm\_client.py} were upgraded with the following safety guardrails:

\begin{guardrailbox}{Guardrail 7: Safe Numpy/Direction Vector Comparison}
NEVER compare Manim direction vectors (like \texttt{UP}, \texttt{DOWN}, \texttt{LEFT}, \texttt{RIGHT}) directly using Python comparison operators (e.g., \texttt{direction == DOWN}). Use \texttt{all(direction == DOWN)} or \texttt{np.array\_equal(direction, DOWN)} to avoid ambiguous numpy evaluation.
\end{guardrailbox}

\begin{guardrailbox}{Guardrail 8: Text Collision Prevention}
NEVER write new text/explanation Mobjects directly on top of existing ones. When updating on-screen explanation/status texts, use \texttt{Transform} or explicitly \texttt{FadeOut} the old text first. Place explanation/status text at the bottom edge using \texttt{.to\_edge(DOWN, buff=0.5)} to avoid spatial collisions with arrays and pointers.
\end{guardrailbox}

\begin{guardrailbox}{Guardrail 9: Screen Boundary Compliance}
Keep horizontal element layouts within screen boundaries (maximum width of 14.2 units). Scale list elements appropriately and keep horizontal spacing buffers small (\texttt{buff=0.1} to \texttt{0.2}) to ensure nothing is clipped at the left/right screen edges.
\end{guardrailbox}

\newpage
\section{Appendix: Full Project Source Code}

Below is the complete, production-ready codebase for the AI Animation Platform.

\subsection{requirements.txt}
\begin{lstlisting}[language=bash]
google-genai
python-dotenv
manim
\end{lstlisting}

\subsection{src/llm\_client.py}
\begin{lstlisting}[language=Python]
import os
import re
import time
from dotenv import load_dotenv
from google import genai
from google.genai import types

# 1. Load the environment variables from the .env file
load_dotenv()

# 2. Extract and check the API key
api_key = os.getenv("GEMINI_API_KEY")
if not api_key:
    raise ValueError(
        "GEMINI_API_KEY is missing from your .env file! "
        "Please ensure your .env file contains: GEMINI_API_KEY=your_key_here"
    )

# 3. Initialize the Google GenAI Client
client = genai.Client(api_key=api_key)

def generate_manim_code(user_instruction: str) -> str:
    """
    Sends a query to Gemini (gemini-2.5-flash) to write Python code
    using the Manim library. Implements exponential backoff retries and strict guardrails.
    """
    system_instruction = (
        "You are an expert mathematical animator using Python's Manim library (Community Edition).\n"
        "Generate a complete, syntactically perfect Python script containing a Manim Scene class representing the animation requested.\n\n"
        "STRICT COMPILATION GUARDRAILS:\n"
        "1. NEVER initialize completely empty text objects like Text(\"\") or Tex(\"\") and then position them using .next_to() or .align_to(). "
        "This triggers a fatal IndexError. Always use at least a space, like Text(\" \"), or a visible placeholder string.\n"
        "2. NEVER use the non-existent '.to_center()' method on any object or VGroup. To center elements, use '.move_to(ORIGIN)' or '.center()'.\n"
        "3. Do not hallucinate or guess methods. Stick to core, standard Manim methods like .shift(), .scale(), .rotate(), .next_to(), .align_to(), and .move_to().\n"
        "4. Ensure all imported modules and custom classes (such as custom Mobjects) are fully defined within this single file.\n"
        "5. Wrap your entire output strictly inside a single code block using standard markdown: ```python ... ```\n"
        "6. Do not write any conversational intro or outro text, explanations, or warnings. Generate ONLY the code block.\n"
        "7. NEVER compare Manim direction vectors (like UP, DOWN, LEFT, RIGHT) directly using Python comparison operators (e.g., `direction == DOWN` or `if direction == UP:`). "
        "Since these are numpy arrays, this raises a ValueError ('The truth value of an array with more than one element is ambiguous'). "
        "Instead, check elements or use `all(direction == DOWN)` or `np.array_equal(direction, DOWN)`.\n"
        "8. NEVER write new text/explanation Mobjects directly on top of existing ones. When updating on-screen explanation/status texts, "
        "use `Transform` or explicitly `FadeOut` the old text before writing/fading in the new one. Place explanation/status text "
        "at the bottom edge of the screen using `.to_edge(DOWN, buff=0.5)` to avoid overlapping with array cells or pointers.\n"
        "9. Keep horizontal element layouts within the screen boundaries. A standard Manim screen has a width of 14.2 units and height of 8.0 units. "
        "When arranging sequences of objects horizontally (like arrays or lists), scale the elements and keep the spacing buffer small "
        "(e.g., `buff=0.1` to `0.2` rather than large values like `0.8`) to prevent objects from being cut off at the left and right edges."
    )

    # Exponential backoff parameters (up to 5 retries: 1s, 2s, 4s, 8s, 16s)
    max_retries = 5
    for attempt in range(max_retries):
        try:
            response = client.models.generate_content(
                model='gemini-2.5-flash',
                contents=user_instruction,
                config=types.GenerateContentConfig(
                    system_instruction=system_instruction,
                    temperature=0.1,  # Keep it highly deterministic
                )
            )
            return response.text
        except Exception as e:
            if attempt == max_retries - 1:
                print(f"\n[Error] API call failed after {max_retries} attempts: {e}")
                return ""
            time.sleep(2 ** attempt)  # Delays of 1s, 2s, 4s, 8s...

def extract_python_code(raw_response: str) -> str:
    """
    Parses out code block delimiters like ```python and ``` to return
    only clean Python code.
    """
    pattern = r"```python\s*(.*?)\s*```"  
    match = re.search(pattern, raw_response, re.DOTALL)  
    if match:  
        return match.group(1).strip()  
    return raw_response.strip()  

def save_code_to_file(code_content: str, filepath: str = "generated_scene.py") -> None:  
    """  
    Overwrites or creates a new file to store the parsed Python script.  
    """  
    try:  
        with open(filepath, "w", encoding="utf-8") as f:  
            f.write(code_content)  
        print(f"\n[Success] Executable Manim animation code successfully saved to: '{filepath}'")  
    except IOError as e:  
        print(f"\n[File Error] Could not write Python code to file: {e}")  

if __name__ == "__main__":  
    print("==========================================================")  
    print(" AI EDUCATIONAL ANIMATION PLATFORM (WEEK 1 LLM BRIDGE) ")  
    print("==========================================================")  

    # Get interactive prompt input directly in your VS Code terminal  
    user_prompt = input("\nEnter your animation concept (e.g. 'Show a blue circle expanding '): ")
    
    if not user_prompt.strip():
        user_prompt = "Create a blue circle in the center that morphs into a yellow square."
        print(f"Using default prompt: '{user_prompt}'")

    print("\nContacting Gemini to generate animation logic...")
    raw_response = generate_manim_code(user_prompt)

    if raw_response:
        print("\n--- Raw Response Received ---")
        print(raw_response)
        print("-----------------------------")

        # Clean and save the response
        clean_code = extract_python_code(raw_response)
        save_code_to_file(clean_code, "generated_scene.py")
    else:
        print("\n[Pipeline Failed] No output was received. Please check your .env API key and connection.")
\end{lstlisting}

\subsection{src/pipeline.py}
\begin{lstlisting}[language=Python]
import os
import re
import subprocess
import sys
from llm_client import generate_manim_code, extract_python_code, save_code_to_file

def find_manim_class_name(code: str) -> str:
    """
    Scans the generated Python code using Regex to find the name of the 
    Manim Scene class so we know which class to render.
    """
    # Look for patterns like "class MyScene(Scene):" or "class MyScene( ZoomedScene ):"
    pattern = r"class\s+(\w+)\s*\(\s*\w*Scene\w*\s*\)\s*:"
    match = re.search(pattern, code)
    if match:
        return match.group(1)
    
    # Fallback to general class finder if "Scene" isn't explicitly in parentheses
    pattern_fallback = r"class\s+(\w+)\s*\("
    match_fallback = re.search(pattern_fallback, code)
    if match_fallback:
        return match_fallback.group(1)
        
    return "ExpandingBlueCircle"  % Default fallback

def run_manim_rendering(filepath: str, class_name: str) -> bool:
    """
    Executes the Manim compiler command as a Python subprocess.
    Equivalent to running: manim -pql filepath class_name in the terminal.
    """
    print(f"\n[Rendering] Compiling class '{class_name}' from '{filepath}'...")
    
    # -p: preview video when finished
    # -ql: low quality (renders lightning fast at 480p)
    command = [
        "manim", 
        "-pql", 
        filepath, 
        class_name
    ]
    
    try:
        # Run the command and stream output directly to VS Code terminal
        result = subprocess.run(command, check=True)
        if result.returncode == 0:
            print("\n[Rendering Success] Animation generated and opened successfully!")
            return True
    except subprocess.CalledProcessError as e:
        print(f"\n[Rendering Error] Manim compilation failed: {e}")
    except FileNotFoundError:
        print(f"\n[System Error] 'manim' command not found. Ensure Manim is installed in your virtual environment.")
    
    return False

def start_pipeline():
    print("==========================================================")
    print("      AI ANIMATION PLATFORM - FULL AUTOMATION PIPELINE    ")
    print("==========================================================")
    
    # 1. Ask user what they want to animate
    user_prompt = input("\nWhat educational animation would you like to build? ")
    if not user_prompt.strip():
        print("Empty input. Running default physics demo...")
        user_prompt = "Animate a projectile motion trajectory with a bouncing ball."
        
    # 2. Get code from Gemini
    raw_response = generate_manim_code(user_prompt)
    if not raw_response:
        print("[Error] Failed to get response from Gemini.")
        return

    # 3. Clean and parse code
    clean_code = extract_python_code(raw_response)
    
    # 4. Save the generated code to a file
    scene_file = "generated_scene.py"
    save_code_to_file(clean_code, scene_file)
    
    # 5. Dynamically extract the name of the Scene class generated by Gemini
    class_name = find_manim_class_name(clean_code)
    print(f"[Parser] Detected generated Scene class name: '{class_name}'")
    
    # 6. Render it automatically!
    run_manim_rendering(scene_file, class_name)

if __name__ == "__main__":
    start_pipeline()
\end{lstlisting}

\end{document}
